\documentclass[a4paper]{article} % Uses the book template for quick start, but scrbook may be better because it is more flexible
% \usepackage[T1]{fontenc} % Include italic fonts
\usepackage[margin=1.5cm, bottom=2cm]{geometry} % Page margins
\usepackage{lettrine}
\usepackage{csquotes}
\usepackage{orcidlink}
\usepackage{amsmath}
\usepackage{graphicx}
\usepackage{minted}
\usepackage{svg}
\usepackage{subcaption}
\usepackage{ccicons}
\usepackage[final]{pdfpages}
\usepackage[inline]{enumitem}

\usepackage{setspace}
\onehalfspacing    % makes the whole document 1.5 spaced
\usepackage{lineno}     % for line numbers

\usepackage[backend=biber, style=numeric]{biblatex}
\addbibresource{erosion.bib}
\usepackage{hyperref}
\hypersetup{colorlinks=true}



\title{Towards stability in erosion models: developing rooted systems in C++ or Rust}
\author{Fee Gevaert \orcidlink{https://orcid.org/0009-0008-2311-6336} \today{}}

\begin{document}
\begin{abstract}
    This essay is aimed to be a guide for people looking to make erosion models.
    In this essay, erosion models, specifically LISEM/OpenLisem are viewed from
    the lens of erosion and erosion control systems. Similar to hillslopes
    affected by erosion, erosion control in the digital domain is complex and
    there are multiple avenues that depend on the local morphology. Unit Tests
    are conceptualized as epistemological roots of an erosion model, ensuring
    soil stability. 
\end{abstract}

\setlist[enumerate]{noitemsep}
\setlist[itemize]{noitemsep}
\linenumbers

\maketitle

\tableofcontents



\section{Introduction}

This essay evolved as a result of trying to run LISEM/OpenLisem on a Linux
system.  Since the recommended manager, MSYS is not available for linux,
dependencies had to be manually collected. The compilation process stranded for
OpenLisem because it depended on a specific commit in a specific branch of QWT
and LISEM similarly depended on a multitude of other softwares. Even though no
successful compilation of either OpenLisem or LISEM was achieved due to their
complexity, I will make some comments with regards to how compilation could have
been easier in section \ref{subs:deps}.

Apart from these comments, I will also propose a workflow for developing erosion
models that should increase maintainability, portability and verifiability of
new or existing codebases. To reify this workflow, I adopted it to port a sample
program \cite{barnesAcceleratingFluvialIncision2019} to Rust. Code is available
on \href{}{Codeberg} and \href{}{Github}.

\begin{enumerate}
    \item Make the algorithm work
    \item Test it on a small (2x2 to 10x10) grid, make those unit tests
    \item verify the tests by hand
    \item Mark this implementation on a separate git branch
    \item Make it run fast, keeping the same tests
\end{enumerate}

Then, when changing unit tests, 

\begin{figure}
    \includesvg[width=\textwidth]{img/catchment_blocks.svg}
    \caption{OpenLisem and its dependencies as a catchment. Possible downstream
    implementations are currently disconnected.}\label{fig:catchment}
\end{figure}

\subsection{OpenLisem Dependency management with Git}\label{subs:deps}

Git has three ways to add third-party repositories:
\href{https://stackoverflow.com/questions/32407634/when-to-use-git-subtree}{Submodules
and Subtrees}. For OpenLisem, I would recommend the use of a Subtree. Sadly, however, I did not get that to work. Therefore, the sad compromise is:


This pulls
in all relevant code from a specific branch, tag or commit, allows for local
changes and is invisible to users of the codebase. Since changing some default
configuration is
\href{https://github.com/vjetten/openlisem/blob/cd0573a24bcb95f3d90c1c36eb1a44b6714123e4/compile_lisem_linux.txt#L35}{advised},
submodules are ruled out.

To add the required branch to the main repository

\printbibliography

\appendix
\end{document}